Machine-learning prediction of enzyme optimal pH has recently been studied directly by Gado et al., who introduced EpHod for predicting enzyme optimum pH from protein sequence \cite{gado2025machine}. Their study evaluated multiple machine-learning methods and reported that models using protein language-model embeddings outperformed other approaches for this task. EpHod is therefore an important reference point for sequence-based optimal-pH prediction, although the present study does not claim to reproduce EpHod because its full training protocol, model architecture, and evaluation setup are not duplicated here.

More broadly, pretrained protein language models have become widely used representations for protein sequence analysis. Early large-scale work showed that unsupervised learning over millions of protein sequences can produce representations containing information relevant to protein structure and function \cite{rives2021biological}. ESM-2 and related evolutionary-scale models further demonstrated that protein language models can support accurate structure-related prediction from sequence alone \cite{lin2023evolutionary}. ProtTrans similarly showed that transformer models trained on large protein-sequence corpora can provide useful embeddings for downstream protein prediction tasks \cite{elnaggar2022prottrans}. These studies motivate the use of frozen protein-language-model embeddings as a feature representation in this work.

A separate line of work emphasizes that evaluation protocols for protein-property prediction must account for sequence relatedness. Hou et al. provide practical guidance for sequence-based protein prediction and specifically warn that homologous proteins shared between training and test sets can inflate measured performance \cite{hou2022ten}. GraphPart formalizes this issue as homology partitioning and argues that closely related biological sequences should be assigned to the same partition when evaluating predictive models \cite{teufel2023graphpart}. This concern motivates the exact-duplicate-safe and cluster-safe split strategies used in the present study.

The provenance and biological validity of pHopt labels also require scrutiny. Zhou and Fang identified fundamental data-quality concerns in the EpHod pHopt resource, including disputed optimal-pH values and protein-identity assignments \cite{zhou2026caution}. Their discussion of accession B5LXD7 is directly relevant here: the local file labels it pH 11.5, whereas the cited experimental literature reports values in the acidic range. We therefore treat such records as audit flags and perform sensitivity analyses rather than interpreting disagreement at an extreme label as model failure.

Sequence clustering tools provide practical mechanisms for constructing these leakage-controlled splits. MMseqs2 enables sensitive and scalable protein sequence searching for large datasets \cite{steinegger2017mmseqs2}, while CD-HIT is a widely used tool for clustering protein or nucleotide sequences at specified identity thresholds \cite{li2006cdhit}. In this project, MMseqs2 was used to generate sequence-identity cluster assignments at 90\%, 70\%, and 50\% thresholds, and all members of each cluster were kept within the same train, validation, or test partition. The lower-identity thresholds provide progressively stricter stress tests of sequence generalization.
